% ============================================================
%  C: Como Programar -- 6ª Edição | Deitel & Deitel
%  Capítulo 8 -- Caracteres e Strings em C
%  FAZENDO A DIFERENÇA (8.43 -- 8.46)
% ============================================================
\documentclass[12pt,a4paper]{article}
\usepackage[utf8]{inputenc}
\usepackage[T1]{fontenc}
\usepackage[brazil]{babel}
\usepackage{geometry}
\geometry{top=2.5cm,bottom=2.5cm,left=3cm,right=3cm}
\usepackage{parskip}\setlength{\parskip}{6pt}\setlength{\parindent}{0pt}
\usepackage{xcolor,tcolorbox,enumitem,listings,fancyhdr,titlesec,hyperref,booktabs,array,amsmath}
\tcbuselibrary{skins,breakable}

\definecolor{deitelblue}{RGB}{0,70,127}
\definecolor{deitelgray}{RGB}{240,242,245}
\definecolor{deitelgreen}{RGB}{0,110,60}
\definecolor{deitellightblue}{RGB}{220,235,250}
\definecolor{deitelred}{RGB}{160,20,20}
\definecolor{deitellorange}{RGB}{200,90,0}
\definecolor{ecogreen}{RGB}{0,120,50}
\definecolor{ecolightgreen}{RGB}{210,240,220}

\newtcolorbox{boaprática}[1][]{colback=deitellightblue,colframe=deitelblue,
  fonttitle=\bfseries\small,title={Boa prática de programação},breakable,#1}
\newtcolorbox{respostabox}[1][]{colback=deitelgray,colframe=deitelblue!60,
  fonttitle=\bfseries\small\color{deitelblue},title={Resposta detalhada},breakable,#1}
\newtcolorbox{enunciadoBox}[1][]{colback=white,colframe=deitelblue,
  fonttitle=\bfseries,title={#1},breakable}
\newtcolorbox{saida}[1][]{colback=black!88,colframe=deitelblue,
  fontupper=\ttfamily\small\color{green!70},
  title={\small\color{white}Saída do programa},breakable,#1}
\newtcolorbox{pesquisa}[1][]{colback=ecolightgreen,colframe=ecogreen,
  fonttitle=\bfseries\small,title={Pesquisa externa realizada},breakable,#1}
\newtcolorbox{atencao}[1][]{colback=yellow!10,colframe=deitellorange,
  fonttitle=\bfseries\small,title={Atenção},breakable,#1}

\setlist[enumerate,1]{label=\textbf{\alph*)},leftmargin=2em}
\setlist[itemize]{leftmargin=2em}

\lstset{
  basicstyle=\ttfamily\small,backgroundcolor=\color{deitelgray},
  frame=single,framerule=0.4pt,rulecolor=\color{deitelblue!40},
  breaklines=true,showstringspaces=false,
  keywordstyle=\color{deitelblue}\bfseries,
  commentstyle=\color{deitelgreen}\itshape,
  stringstyle=\color{deitelred},
  language=C,numbers=left,numberstyle=\tiny\color{gray},
  xleftmargin=18pt,xrightmargin=5pt,stepnumber=1,
}

\pagestyle{fancy}\fancyhf{}
\fancyhead[L]{\small\color{deitelblue}\textbf{C: Como Programar -- 6ª ed. | Deitel \& Deitel}}
\fancyhead[R]{\small\color{deitelblue}Cap. 8 -- Fazendo a Diferença}
\fancyfoot[C]{\small\thepage}
\renewcommand{\headrulewidth}{0.4pt}

\titleformat{\section}[block]{\large\bfseries\color{deitelblue}}{\thesection.}{0.5em}{}[\titlerule]
\titleformat{\subsection}[block]{\normalsize\bfseries\color{deitelblue!80}}{}{}{}

\hypersetup{colorlinks=true,linkcolor=deitelblue,urlcolor=ecogreen}

\begin{document}

\begin{titlepage}
  \centering\vspace*{2cm}
  {\color{deitelblue}\rule{\linewidth}{2pt}}\\[0.4cm]
  {\huge\bfseries\color{deitelblue}C: Como Programar}\\[0.2cm]
  {\large\color{deitelblue}6ª Edição $\cdot$ Paul Deitel \& Harvey Deitel}\\[0.2cm]
  {\color{deitelblue}\rule{\linewidth}{2pt}}\\[1cm]
  {\LARGE\bfseries Capítulo 8}\\[0.3cm]
  {\Large\color{deitelblue!80}Caracteres e Strings em C}\\[0.8cm]
  \begin{tcolorbox}[colback=ecolightgreen,colframe=ecogreen,width=0.85\linewidth]
    \centering{\large\bfseries\color{ecogreen}Fazendo a Diferença}\\[0.3cm]
    {\normalsize Exercícios 8.43 a 8.46\\[0.2cm]
    Receitas saudáveis · Detector de spam\\
    Linguagem SMS · Neutralidade de gênero}
  \end{tcolorbox}
  \vfill{\small Abordagem incremental Deitel}
\end{titlepage}

\tableofcontents\newpage

% ============================================================
\section{Exercício 8.43 -- Cozinhando com Ingredientes Saudáveis}
% ============================================================

\begin{enunciadoBox}[Enunciado 8.43]
Escreva um programa que leia uma receita do usuário e sugira substitutos mais
saudáveis para ingredientes. O programa deve considerar problemas como colesterol
alto, perda de peso e alergias. Use as substituições típicas da tabela do livro
(Figura 8.40).
\end{enunciadoBox}

\subsection*{Pesquisa: Substituições de Ingredientes}

\begin{pesquisa}
Conforme a tabela da Figura 8.40 do livro (com fontes em sites de nutrição):
\begin{itemize}[noitemsep]
  \item 1 xícara de creme de leite $\to$ 1 xícara de iogurte
  \item 1 xícara de açúcar $\to$ 1/2 xícara de mel ou 1/4 de adoçante
  \item 1 xícara de manteiga $\to$ 1 xícara de margarina/iogurte
  \item 1 ovo $\to$ 2 colheres de amido + 2 claras (ou banana amassada)
  \item Farinha branca $\to$ farinha integral ou de arroz
  \item 1 xícara de maionese $\to$ 1 xícara de cottage
\end{itemize}

A lógica do programa: percorrer a receita, identificar ingredientes-alvo via
\texttt{strstr}, e substituí-los conforme o perfil de saúde do usuário.
\end{pesquisa}

\subsection*{Solução em C}

\begin{respostabox}
\begin{lstlisting}
/* Ex8.43: receita_saudavel.c
   Sugere substituicoes mais saudaveis em receitas */
#include <stdio.h>
#include <string.h>
#include <ctype.h>
#define MAX_RECEITA 2000

/* Estrutura de substituicao */
typedef struct {
    const char *original;
    const char *substituto;
    const char *motivo;  /* "colesterol", "perda peso", "alergia" */
} Substituicao;

void sugerir( const char *receita, const Substituicao subs[], int n );

int main( void )
{
    Substituicao subs[] = {
        { "creme de leite", "iogurte natural",           "colesterol/peso" },
        { "manteiga",       "azeite ou margarina light", "colesterol" },
        { "acucar",         "mel ou adocante",            "peso/diabetes" },
        { "maionese",       "queijo cottage",             "colesterol" },
        { "ovo",            "2 claras + 2 colheres amido","colesterol" },
        { "farinha branca", "farinha integral",            "perda peso"   },
        { "leite",          "leite desnatado/soja",        "colesterol" }
    };
    int n = sizeof( subs ) / sizeof( subs[ 0 ] );

    char receita[ MAX_RECEITA ];

    printf( "===== CONVERSOR DE RECEITAS SAUDAVEIS =====\n\n" );
    printf( "ATENCAO: Sempre consulte um medico antes de fazer\n" );
    printf( "         mudancas significativas em sua dieta.\n\n" );

    printf( "Digite sua receita (digite END em uma linha sozinha):\n" );

    /* Le linhas ate encontrar "END" */
    receita[ 0 ] = '\0';
    char linha[ 200 ];
    while ( fgets( linha, sizeof( linha ), stdin ) != NULL ) {
        if ( strncmp( linha, "END", 3 ) == 0 ) break;
        strcat( receita, linha );
    }

    sugerir( receita, subs, n );

    return 0;
}

void sugerir( const char *receita, const Substituicao subs[], int n )
{
    int i;
    int encontrou_algo = 0;

    printf( "\n=========== SUGESTOES ===========\n" );

    for ( i = 0; i < n; ++i ) {
        if ( strstr( receita, subs[ i ].original ) != NULL ) {
            printf( "- Substitua \"%s\" por \"%s\"\n",
                    subs[ i ].original, subs[ i ].substituto );
            printf( "  (beneficio: %s)\n", subs[ i ].motivo );
            encontrou_algo = 1;
        }
    }

    if ( !encontrou_algo ) {
        printf( "Nenhuma substituicao aplicavel encontrada.\n" );
        printf( "Sua receita parece estar com ingredientes ja saudaveis!\n" );
    }
}
\end{lstlisting}
\end{respostabox}

\newpage

% ============================================================
\section{Exercício 8.44 -- Analisador de Spam}
% ============================================================

\begin{enunciadoBox}[Enunciado 8.44]
Pesquise palavras comuns em spam e crie uma lista com 30. Escreva um programa que
leia uma mensagem (sem ultrapassar o limite do array), conte ocorrências dessas
palavras-chave, e avalie a probabilidade de a mensagem ser spam.
\end{enunciadoBox}

\subsection*{Solução em C}

\begin{respostabox}
\begin{lstlisting}
/* Ex8.44: analisador_spam.c */
#include <stdio.h>
#include <string.h>
#include <ctype.h>
#define MAX_MSG 5000
#define NUM_PALAVRAS 30

int main( void )
{
    const char *palavrasSpam[ NUM_PALAVRAS ] = {
        "gratis",      "promocao",    "ganhe",       "premio",
        "credito",     "oferta",      "dinheiro",    "urgente",
        "exclusivo",   "limitado",    "100%",        "garantido",
        "renda",       "milionario",  "investimento", "lucro",
        "viagra",      "remedio",     "farmacia",    "barato",
        "desconto",    "imperdivel",  "milagre",     "secreto",
        "rapido",      "sucesso",     "saude",       "emagrecer",
        "clique aqui", "ganhador"
    };

    char msg[ MAX_MSG ];
    char msgLower[ MAX_MSG ];
    int  pontos = 0;
    int  i, j;
    char linha[ 200 ];

    printf( "Digite a mensagem de e-mail (END para terminar):\n" );

    msg[ 0 ] = '\0';
    while ( fgets( linha, sizeof( linha ), stdin ) != NULL ) {
        if ( strncmp( linha, "END", 3 ) == 0 ) break;

        /* Garante que nao excede o tamanho */
        if ( strlen( msg ) + strlen( linha ) < MAX_MSG - 1 ) {
            strcat( msg, linha );
        }
    }

    /* Converte para minusculas para busca case-insensitive */
    for ( i = 0; msg[ i ] != '\0'; ++i ) {
        msgLower[ i ] = tolower( ( unsigned char )msg[ i ] );
    }
    msgLower[ i ] = '\0';

    /* Conta ocorrencias de cada palavra-chave */
    printf( "\n=== ANALISE ===\n" );
    for ( i = 0; i < NUM_PALAVRAS; ++i ) {
        int conta = 0;
        char *p = msgLower;
        while ( ( p = strstr( p, palavrasSpam[ i ] ) ) != NULL ) {
            ++conta;
            ++p;
        }
        if ( conta > 0 ) {
            printf( "  '%s' aparece %d vez(es)\n", palavrasSpam[ i ], conta );
            pontos += conta;
        }
    }

    /* Avaliacao */
    printf( "\nNivel de spam: %d pontos\n", pontos );
    if ( pontos == 0 ) {
        printf( "Provavelmente NAO e spam.\n" );
    }
    else if ( pontos <= 2 ) {
        printf( "Possivelmente legitimo.\n" );
    }
    else if ( pontos <= 5 ) {
        printf( "SUSPEITO: pode ser spam.\n" );
    }
    else {
        printf( "MUITO PROVAVELMENTE SPAM!\n" );
    }

    return 0;
}
\end{lstlisting}
\end{respostabox}

\subsection*{Reflexão}

\begin{boaprática}
Filtros de spam reais usam técnicas muito mais sofisticadas:
\begin{itemize}[noitemsep]
  \item \textbf{Bayes ingênuo:} cálculo probabilístico baseado em palavras e suas frequências em corpus de spam vs.\ ham.
  \item \textbf{Análise de cabeçalhos:} verificação de remetente, SPF, DKIM.
  \item \textbf{Listas de bloqueio (blacklists):} domínios e IPs conhecidos por enviar spam.
  \item \textbf{Aprendizado de máquina:} redes neurais treinadas em milhões de mensagens.
\end{itemize}
O exercício é uma porta de entrada para esse domínio fascinante, em que estatística
e linguística computacional encontram a engenharia de software.
\end{boaprática}

\newpage

% ============================================================
\section{Exercício 8.45 -- Linguagem SMS}
% ============================================================

\begin{enunciadoBox}[Enunciado 8.45]
Escreva um programa que traduza mensagens em \textit{linguagem SMS} (abreviações
de mensagens de texto) para português padrão e vice-versa. Lembre-se de que
algumas abreviações podem ter múltiplos significados (ambiguidade).
\end{enunciadoBox}

\subsection*{Solução em C}

\begin{respostabox}
\begin{lstlisting}
/* Ex8.45: tradutor_sms.c */
#include <stdio.h>
#include <string.h>
#include <ctype.h>
#define MAX_MSG 1000

typedef struct {
    const char *sms;
    const char *portugues;
} Abreviacao;

void substituirTodas( char *texto, const char *de, const char *para );

int main( void )
{
    Abreviacao tabela[] = {
        { "vc",  "voce" },        { "tb",  "tambem" },
        { "blz", "beleza" },      { "pq",  "porque" },
        { "tbm", "tambem" },      { "qq",  "qualquer" },
        { "qdo", "quando" },      { "msm", "mesmo" },
        { "td",  "tudo" },        { "qto", "quanto" },
        { "obg", "obrigado" },    { "vlw", "valeu" },
        { "abs", "abracos" },     { "bjs", "beijos" },
        { "emo", "em minha opiniao" }
    };
    int n = sizeof( tabela ) / sizeof( tabela[ 0 ] );

    int opcao;
    char texto[ MAX_MSG ];
    int  i;

    printf( "1 - SMS -> Portugues\n2 - Portugues -> SMS\nEscolha: " );
    scanf( "%d", &opcao );
    getchar();  /* consome \n */

    printf( "Digite a mensagem: " );
    fgets( texto, MAX_MSG, stdin );

    if ( opcao == 1 ) {
        for ( i = 0; i < n; ++i ) {
            substituirTodas( texto, tabela[ i ].sms, tabela[ i ].portugues );
        }
    }
    else {
        for ( i = 0; i < n; ++i ) {
            substituirTodas( texto, tabela[ i ].portugues, tabela[ i ].sms );
        }
    }

    printf( "\nResultado: %s\n", texto );
    return 0;
}

void substituirTodas( char *texto, const char *de, const char *para )
{
    char buffer[ MAX_MSG ];
    char *p;
    int  tamDe = strlen( de );
    int  tamPara = strlen( para );

    while ( ( p = strstr( texto, de ) ) != NULL ) {
        /* Copia ate o ponto encontrado */
        int prefixo = p - texto;
        strncpy( buffer, texto, prefixo );
        buffer[ prefixo ] = '\0';

        /* Concatena substituicao */
        strcat( buffer, para );

        /* Concatena resto apos o termo encontrado */
        strcat( buffer, p + tamDe );

        /* Copia de volta */
        strcpy( texto, buffer );
    }
}
\end{lstlisting}

\textbf{Problema da ambiguidade:} ``EMO'' pode significar ``em minha opinião'',
``Empresa Marítima Ocidental'', ou um gênero musical. Sistemas reais usam contexto
para resolver: vizinhança de outras palavras, modelos estatísticos de linguagem,
ou perguntar ao usuário. Esse é um problema clássico de \textit{desambiguação
lexical}.
\end{respostabox}

\newpage

% ============================================================
\section{Exercício 8.46 -- Neutralidade de Gênero}
% ============================================================

\begin{enunciadoBox}[Enunciado 8.46]
Crie um programa que leia um parágrafo de texto e substitua palavras de gênero
específico por equivalentes sem gênero. Apresente o texto resultante.
\end{enunciadoBox}

\subsection*{Solução em C}

\begin{respostabox}
\begin{lstlisting}
/* Ex8.46: neutralizar_genero.c */
#include <stdio.h>
#include <string.h>
#include <ctype.h>
#define MAX_TEXTO 5000

typedef struct {
    const char *generificado;
    const char *neutro;
} Substituicao;

void substituirTodas( char *texto, const char *de, const char *para );

int main( void )
{
    /* Tabela de substituicoes para neutralidade de genero
       (apenas alguns exemplos -- expandir conforme a necessidade) */
    Substituicao subs[] = {
        { "porta-voz",   "representante"      },
        { "presidente",  "presidencia"        },
        { "homem-feito", "industrial"         },
        { "homens",      "pessoas"            },
        { "homem",       "pessoa"             },
        { "policial",    "policial"           },  /* ja neutro */
        { "bombeiro",    "agente do corpo de bombeiros" },
        { "garcom",      "atendente"          },
        { "carteiro",    "agente postal"      },
        { "merceeiro",   "comerciante"        },
        { "marinheiro",  "tripulante"         },
        { "doutora",     "doutor(a)"          }
    };
    int n = sizeof( subs ) / sizeof( subs[ 0 ] );

    char texto[ MAX_TEXTO ];
    char linha[ 500 ];
    int i;

    printf( "Digite o paragrafo (END para terminar):\n" );

    texto[ 0 ] = '\0';
    while ( fgets( linha, sizeof( linha ), stdin ) != NULL ) {
        if ( strncmp( linha, "END", 3 ) == 0 ) break;
        if ( strlen( texto ) + strlen( linha ) < MAX_TEXTO - 1 ) {
            strcat( texto, linha );
        }
    }

    /* Aplica todas as substituicoes */
    for ( i = 0; i < n; ++i ) {
        substituirTodas( texto, subs[ i ].generificado, subs[ i ].neutro );
    }

    printf( "\n=== Texto neutralizado ===\n%s\n", texto );

    return 0;
}

void substituirTodas( char *texto, const char *de, const char *para )
{
    char buffer[ MAX_TEXTO ];
    char *p;
    int  tamDe = strlen( de );

    while ( ( p = strstr( texto, de ) ) != NULL ) {
        int prefixo = p - texto;
        strncpy( buffer, texto, prefixo );
        buffer[ prefixo ] = '\0';
        strcat( buffer, para );
        strcat( buffer, p + tamDe );
        strcpy( texto, buffer );
    }
}
\end{lstlisting}
\end{respostabox}

\subsection*{Reflexão}

\begin{atencao}
A neutralidade linguística em português é um tema mais complexo do que em inglês
por causa da \textit{flexão de gênero em substantivos e adjetivos}. Soluções
profissionais consideram:

\begin{itemize}[noitemsep]
  \item \textbf{Reescrita estrutural:} ``o aluno deve fazer'' $\to$ ``é necessário fazer''
  \item \textbf{Formas duplas:} ``aluno/aluna'', ``aluno(a)''
  \item \textbf{Linguagem inclusiva moderna:} ``alunes'', ``alunx'' (controverso)
  \item \textbf{Coletivos:} ``estudantes'', ``pessoa interessada''
\end{itemize}

A substituição automática é apenas o primeiro passo. Sistemas avançados usam
processamento de linguagem natural (PLN) para entender contexto e propor
reescritas mais naturais.
\end{atencao}

\bigskip

\begin{boaprática}
\textbf{Padrão comum nos 4 exercícios desta seção:} todos envolvem
\textit{substituição contextual} em texto, e implementam o algoritmo central
\texttt{substituirTodas} (busca + substituição). Reusabilidade!

Este é o princípio do \textbf{DRY} (\textit{Don't Repeat Yourself}): identificar
o padrão e implementar uma vez. O Cap.~5 (funções) e o Cap.~7 (ponteiros) deram
a você as ferramentas necessárias para fazê-lo elegantemente.
\end{boaprática}

\newpage

% ============================================================
\section*{Reflexão Final do Capítulo 8}

\begin{tcolorbox}[colback=deitellightblue,colframe=deitelblue,
  title={\bfseries A computação a serviço da linguagem humana}]

O Capítulo~8 nos leva a um território fascinante: a interseção entre programação
e linguagem natural. Os exercícios ``Fazendo a Diferença'' demonstram aplicações
de manipulação de strings que tocam questões sociais profundas:

\begin{itemize}[noitemsep]
  \item \textbf{Saúde pública} (8.43) -- ajudar pessoas a comer melhor.
  \item \textbf{Segurança digital} (8.44) -- proteger usuários de e-mails maliciosos.
  \item \textbf{Comunicação} (8.45) -- traduzir entre dialetos da era digital.
  \item \textbf{Inclusão social} (8.46) -- promover linguagem mais inclusiva.
\end{itemize}

Em todos eles, o ponto central é o mesmo: \textbf{texto é uma forma de dados
extremamente importante}, e dominar suas manipulações abre portas para
aplicações com impacto real na sociedade.

\textbf{Aonde isso leva?} O processamento de linguagem natural (PLN) é hoje
uma das áreas mais quentes da computação. Modelos de linguagem como GPT, BERT
e Claude são, no fundo, sistemas extraordinariamente sofisticados de manipulação
de strings -- com bilhões de parâmetros, mas baseados em ideias fundamentais que
você acabou de praticar: tokenização, busca, substituição, classificação.

\end{tcolorbox}

\begin{boaprática}
\textbf{Próximos passos:} no Capítulo~9, estudaremos formatação detalhada de
E/S; no Capítulo~10, estruturas (\texttt{struct}) que permitem agrupar dados
heterogêneos; no Capítulo~11, manipulação de arquivos (essencial para
processamento de textos do mundo real); no Capítulo~12, estruturas de dados
encadeadas (que suportam dicionários, tabelas hash, parsing de árvores
sintáticas, etc.). O que você acabou de aprender é a base de tudo.
\end{boaprática}

\end{document}
